Después de la optimización de hiperparámetros y parámetros de cada modelo, y de la toma de decisión de arquitectura en el caso de LSTM, se llegó a los modelos representativos de la Tabla~\ref{tab:modelos_seleccionados}:
\begin{table}[htbp]
\centering
\small
\begin{tabular}{lp{9.5cm}}
\toprule
Modelo & Arquitectura / Parámetros\\
\midrule
LSTM & (*) \\
Prophet & changepoint prior scale: 0.016, seasonality prior scale: 750.108 \\
LGBM & colsample bytree: 0.64, learning rate: 0.10, max depth: 6, n estimators: 63, num leaves: 67 \\
XGBoost & colsample bytree: 0.37, learning rate: 0.02,  max depth: 10, n estimators: 190 \\
ARIMA & d: 0, p: 29, q: 24, trend: c \\
\bottomrule
\end{tabular}
\caption{Modelos seleccionados y su composición}
\label{tab:modelos_seleccionados}
\end{table}


(*) La arquitectura y los hiperparámetros del modelo LSTM elegido son:

\subsubsection*{Arquitectura de la LSTM}

- Capa de entrada con dimensiones de entrada (batch size, time steps, features) igual a (None, 10, 17).

- Capa LSTM con 2700 unidades, activación 'relu', función de activación recurrente 'sigmoid', y una tasa de dropout del 30\%. Esta capa regresa secuencias y no estados.

- Capa de dropout con una tasa de dropout del 30\%, aplicada a la salida de la capa LSTM.

- Capa LSTM con 800 unidades, activación 'relu', función de activación recurrente 'sigmoid', y una tasa de dropout del 30\%. Esta capa no retorna secuencias, solo el estado final.

- Capa de dropout con una tasa de dropout del 30\%, aplicada a la salida de la segunda capa LSTM.

- Capa densa con 1 unidad y activación 'relu'.

\subsubsection*{Hiperparámetros de la LSTM}

- Unidades en la primera capa LSTM: 2700.

- Unidades en la segunda capa LSTM: 800.

- Tasa de dropout: 30\%.

- Función de activación en las capas LSTM: 'relu'.

- Función de activación recurrente en las capas LSTM: 'sigmoid'.

- Se utiliza la inicialización de pesos GlorotUniform para las capas LSTM y Dense.

- Se utiliza la inicialización de pesos Orthogonal para los pesos recurrentes de las capas LSTM.

- Las capas LSTM no son stateful y no se desenrollan en el tiempo.

- Se utiliza la implementación 2 para las capas LSTM.

- La capa densa tiene 1 unidad de salida.\\

Los resultados de las métricas de cada modelo, tal como los reportó el borrador, se resumen en la Tabla~\ref{tab:metricas_borrador}; las gráficas de valores predichos contra valores reales de cada modelo están en el Apéndice~\ref{ap:figuras} (Figuras \ref{fig:pred_arima} a \ref{fig:pred_ensemble_ambas}).

\begin{table}[htbp]
\centering
\small
\begin{tabular}{lccccc}
\toprule
Modelo & Err.\ temp.\ 2020-21 & Err.\ temp.\ 2021-22 & Err.\ ambas & RMSE & MAE \\
\midrule
XGBoost + LSTM & 8.7\% & 92\% & 12.9\% & 527.3 & 369.0 \\
XGBoost & 12\% & 16.3\% & 14.5\% & 450.7 & 312.6 \\
LGBM & 36\% & 36\% & 36.4\% & 494.9 & 380.2 \\
\midrule
LSTM & --- & --- & 61\% & 183.1 & 104.9 \\
Prophet & 164\% & 1737\% & 761\% & 135.0 & 105.8 \\
ARIMA & 338\% & 142\% & 142\% & 467.5 & 344.1 \\
\bottomrule
\end{tabular}
\caption{Resultados de la primera iteración (borrador de 2023) por modelo: error en la suma de
temporada (2020-2021, 2021-2022 y ambas) y RMSE/MAE diarios en kg. \emph{Nota}: las cifras se
transcriben tal cual del borrador, cuyo proceso de datos dependía de fuentes hoy inaccesibles y
no es reproducible; contienen inconsistencias internas menores de transcripción (algunos errores
porcentuales no cuadran exactamente con las sumas que los acompañaban), por lo que deben leerse
como el registro histórico del punto de partida y no como cifras verificadas. La re-corrida
verificable de estos mismos modelos sobre el conjunto unido 2017-2026 está en la
Tabla~\ref{tab:extension_comparativa}.}
\label{tab:metricas_borrador}
\end{table}


Se optó por ensamblar XGBoost, el mejor modelo para predecir la tendencia, junto con LSTM: primero se hace la predicción con XGBoost y ese vector de y se introduce en LSTM como un regresor más. Los resultados del ensamble son los mejores, con un error de captura del 8.7\% para la primera temporada y un error global del 12.9\%. A pesar de esto, los resultados para la segunda temporada no son nada favorables. 

Además, los RMSE y MAE son elevados, lo cual implica que las predicciones diarias son poco precisas, a pesar de que la suma de toda la temporada se predice con poco error. Esto puede significar un reto para confiar en el modelo, pero es consecuencia de la falta de granularidad en los datos de entrenamiento.

En general, todos los modelos sobreestimaron la pesca de la segunda temporada debido a que podría tratarse de un caso atípico. Si comparamos esta temporada con el resto, el total de pesca es menos de la mitad. Por esta razón se decidió optar por predecir la temporada 2020-2021 también, pues es similar a las demás.

No obstante, debemos poner atención a la causa de la reducción de la pesca en la temporada 2021-2022. Inicialmente se planteó la hipótesis de un efecto económico externo (el último año de la pandemia de Covid-19); sin embargo, el índice de olas de calor marinas de la Sección~\ref{sec:datos_mhw} atribuye la caída a un mecanismo ambiental: las olas de calor marinas que afectaron a Baja California entre 2019 y 2021 \citep{villasenor2024}. Esto explica por qué los modelos que solo usan el historial de capturas —o las variables oceanográficas sin un índice de anomalía térmica de largo plazo— sobreestiman sistemáticamente esa temporada.

El paso siguiente que este planteamiento dejaba señalado ---conseguir datos de otras regiones y de otras especies, replicar el procedimiento y verificar si los modelos elegidos sostienen sus resultados--- es exactamente el que emprenden las secciones que siguen.
